% latex/catalogue.tex — printable catalogue shell (BBL-208 redesign, v1.2.0).
% \input{}s two GENERATED files (never hand-edited):
%   catalogue_meta.tex  — written by scripts/toledo_build.py (build_catalogue_meta):
%                         \ToledoVersion / \ToledoLicense, read from CITATION.cff.
%   catalogue_body.pdf.tex — derived by scripts/latex_pdf_safe.py from the GENERATED
%                         catalogue_body.tex (written by scripts/toledo_build.py's
%                         build_catalogue_body): Thai/Cyrillic runs -> a disclosed
%                         placeholder, and any display-math (\[ \]) block that fails
%                         to compile standalone -> the same wrapped-monospace fallback
%                         used for ascii-math (also disclosed). The registries, JSON entries,
%                         docs site and vault all carry the exact source text
%                         unmodified; only this print artifact ever substitutes.
%
% Statements are copied verbatim from their sources (BBL-172, latest formulation
% wins) and carry literal Unicode this machine's pdflatex cannot typeset natively
% (math alphanumerics, arrows/relations, Thai, Cyrillic) — see
% latex/unicode_pdf_fallback.sty (\newunicodechar) and scripts/latex_pdf_safe.py.
\documentclass[10pt,a4paper]{article}
\usepackage[T1]{fontenc}
\usepackage{lmodern}
\usepackage{textcomp}
\usepackage{amsfonts}
\usepackage{amsmath}
\usepackage{pifont}
\usepackage{newunicodechar}
\usepackage[margin=2.2cm]{geometry}
\usepackage{booktabs}
\usepackage{makeidx}
\usepackage{xcolor}
\usepackage[colorlinks=true,linkcolor=blue!40!black,urlcolor=blue!40!black,citecolor=blue!40!black]{hyperref}
\usepackage{toledo}
\usepackage{unicode_pdf_fallback}

\makeindex
\allowdisplaybreaks
\setcounter{tocdepth}{1}   % Parts (-1) + root Sections (1) only — individual
                           % entries are \subsubsection* (starred, never added
                           % to the ToC), and status_counts_block's \subsection*
                           % blocks in End matter are starred too.
\setcounter{secnumdepth}{1}

\title{Toledo}
\date{}
\author{}

\input{catalogue_meta.tex}

\begin{document}

\begin{titlepage}
\centering
\vspace*{2cm}
{\fontsize{34}{40}\selectfont\bfseries Toledo\par}
\vspace{0.6cm}
{\Large Equation Library of the Human--AI Readout Programme\par}
\vspace{2cm}
{\large Version \ToledoVersion\par}
\vspace{0.4cm}
Concept DOI: \href{https://doi.org/10.5281/zenodo.22537318}{10.5281/zenodo.22537318}\par
\vspace{0.4cm}
Generated \today\par
\vspace{0.4cm}
Licence: \ToledoLicense\par
\vfill
\begin{minipage}{0.86\textwidth}
\small\noindent
\textbf{How to read this catalogue.} Every object carries a \emph{code}. A
\textbf{root} code is the source's own identifier, copied verbatim and never
renumbered (an Appendix~C stream id such as \texttt{EQ-008}, a named result such
as \texttt{weld} or \texttt{MQ.08}, or a founder-ruled root-registry extension such
as \texttt{Theta}/\texttt{CMC}). A \textbf{reading} code has the shape
\texttt{<root>/<D>.<nn>.v<k>}: \texttt{<D>} is one of eight domain letters
(E~epistemic, H~human--AI, S~social, W~world-system, M~method, P~physics,
C~chemistry, B~biology/health) naming which domain-lens the root is read through;
\texttt{<nn>} is that reading's sequence number under its root; \texttt{v<k>} is
the revision of its statement (the latest revision is always shown). Every entry
also states its \emph{tier} (an honest evidence label, copied from the source and
never raised): \texttt{Th\_coqc} (machine-checked in Coq) and
\texttt{finite\_diagnostic}/\texttt{Dr} (checked by a weaker, disclosed method) are
the tiers with the strongest evidence behind them; \texttt{Open} discloses an
acknowledged open question; \texttt{Definition}, \texttt{Ax} (a named, disclosed
axiom) and \texttt{RETRACTED} are categorical rather than evidence-graded;
\texttt{untagged} means the source attached no tier marker at all. Separately,
\emph{coq\_status} states what machine-checked Coq artifact (if any) backs the
entry today: \texttt{closed} (a Toledo-native file proves it, ``Closed under the
global context''), \texttt{wrapped\_related} (a Toledo-native wrapper aliases an
imported identifier that only reads or specialises the statement),
\texttt{mapped\_not\_wrapped} (an evidence-backed match in the Coq import exists
but no Toledo-native wrapper file yet), \texttt{definition} (a typed
Definition/Record, no theorem), \texttt{open\_prop} (an Open/Dr hypothesis stated
as an unproved Prop), \texttt{not\_formalisable} (no formal content in the
source), or \texttt{none}/\texttt{not\_yet\_formalised}. A statement typeset as
display mathematics is the source's own LaTeX transcription; a statement shown in
a monospaced block is the source's own ASCII-math or Coq text, wrapped for print
but otherwise unedited; a plain paragraph is prose the source itself states has no
operator. \emph{Parents} lists the code(s) this entry is derived from or reads;
\emph{occurrences} lists where it was found in the corpus.
\end{minipage}
\vspace{1cm}
\end{titlepage}

\clearpage
\tableofcontents
\clearpage

\input{catalogue_body.pdf.tex}

\end{document}
